\section{Formal Notation}
\label{app:formal-notation}

This appendix fixes the mathematical vocabulary used in this paper. Graphs are
finite, and node identifiers are implementation handles rather than semantic
attributes. Literal equality is denoted by \(=\); equivalence under a declared
isomorphism or relabelling relation is denoted by \(\equiv\).

\subsection{Summary of symbols}

\begin{table}[ht]
  \centering
  \caption{Core notation for Abstract Graph representations.}
  \label{tab:formal-notation}
  \begin{tabular}{@{}ll@{}}
    \toprule
    Symbol & Meaning \\
    \midrule
    \(G=(V,E,X_V,X_E,\tau_G)\) & Attributed base graph and graph kind \\
    \(G[U]\), \(G\langle F\rangle\) & Node- and edge-induced subgraphs \\
    \(\mathcal{S}(G)\) & Permitted mapped-subgraph space of \(G\) \\
    \(I=(V_I,E_I,X_I,R_I)\) & Interpretation graph and relation semantics \\
    \(\mu:V_I\to\mathcal{S}(G)\) & Interpretation-node mapping \\
    \(A=(G,I,\mu)\) & Abstract Graph representation \\
    \(\Gamma=(\ell,a,r,\theta)\) & Execution configuration \\
    \(s_1\equiv_q s_2\) & Attributed-subgraph equivalence \\
    \(c_q(s)\) & Canonical structural certificate \\
    \(\phi_{q,b}(s)\) & Structural feature identity at width \(b\) \\
    \(\pi_G(y)\), \(\delta(y)\) & Extensional and derivation provenance \\
    \(o\), \(P\) & Operator and operator program \\
    \bottomrule
  \end{tabular}
\end{table}

\subsection{Base graphs and mapped subgraphs}

\begin{definition}[Base graph]
An attributed base graph is
\begin{equation}
  G=(V,E,X_V,X_E,\tau_G),
\end{equation}
where \(V\) and \(E\) are its node and edge sets, \(X_V\) and \(X_E\) are
node- and edge-attribute maps, and \(\tau_G\) records the graph kind, including
directedness. This paper considers finite simple directed and undirected
graphs. Multigraphs and hypergraphs are outside its scope.
\end{definition}

For \(U\subseteq V\), the node-induced subgraph is
\begin{equation}
  G[U]=\bigl(U,E\cap(U\times U)\bigr).
\end{equation}
For \(F\subseteq E\), \(G\langle F\rangle\) denotes the edge-induced subgraph
containing \(F\) and its incident nodes. Semantic attributes are restricted to
the selected nodes and edges.

\begin{definition}[Mapped-subgraph space]
The permitted mapped-subgraph space \(\mathcal{S}(G)\) contains objects
\begin{equation}
  s=(V_s,E_s,X_V|_{V_s},X_E|_{E_s}),
  \qquad V_s\subseteq V,\quad E_s\subseteq E.
\end{equation}
Node- and edge-induced mapped subgraphs are allowed when the producing operator
declares its materialisation rule. Mapped subgraphs may overlap, equal the
whole base graph, and be disconnected. Empty mappings require an explicit
operator contract and are suppressed by default in reference decompositions.
\end{definition}

\subsection{Interpretation graphs and Abstract Graphs}

\begin{definition}[Interpretation graph]
An interpretation graph is
\begin{equation}
  I=(V_I,E_I,X_I,R_I),
\end{equation}
where every \(u\in V_I\) denotes a mapped subgraph, \(X_I\) contains derived
attributes, and \(R_I\) declares interpretation-edge semantics. An
interpretation edge may encode overlap, derivation, compatibility, or another
named relation; it need not denote base-graph adjacency.
\end{definition}

\begin{definition}[Mapping and Abstract Graph]
The total mapping
\begin{equation}
  \mu:V_I\longrightarrow\mathcal{S}(G)
\end{equation}
associates every interpretation node with exactly one mapped subgraph. It need
not be injective or cover all of \(G\). An Abstract Graph is
\begin{equation}
  A=(G,I,\mu).
\end{equation}
\end{definition}

Construction functions remain separate from the semantic triple. When
relevant, the execution configuration is
\begin{equation}
  \Gamma=(\ell,a,r,\theta),
\end{equation}
where \(\ell\) assigns labels or identities, \(a\) computes derived
attributes, \(r\) constructs interpretation relations, and \(\theta\) contains
their parameters. An implementation state may be written \((A,\Gamma)\).

\subsection{Equivalence and structural identity}

A permitted relabelling is a bijection \(\rho:V\to V'\) preserving graph kind,
adjacency or arc direction, and all attributes selected as semantic. Let \(q\)
denote that attribute projection.

\begin{definition}[Mapped-subgraph equivalence]
Two mapped subgraphs are equivalent, written \(s_1\equiv_q s_2\), if an
isomorphism between them preserves directedness and all node and edge
attributes selected by \(q\).
\end{definition}

Let \(c_q(s)\) be a canonical certificate and \(h_b\) a bounded hash of width
\(b\). The structural feature identity is
\begin{equation}
  \phi_{q,b}(s)=h_b\bigl(c_q(s)\bigr).
\end{equation}
The intended certificate property is
\begin{equation}
  s_1\equiv_q s_2 \Longrightarrow c_q(s_1)=c_q(s_2).
\end{equation}
The converse depends on the canonicalisation method. Equality of bounded
feature identities does not prove equivalence because hashing may introduce
collisions. Identity experiments therefore report \(q\), \(b\), the
certificate and hash algorithms, and the collision policy.

\begin{definition}[Abstract Graph equivalence]
Let \(A=(G,I,\mu)\) and \(A'=(G',I',\mu')\). We write
\begin{equation}
  A\equiv_{q,R_I}A'
\end{equation}
if there exist graph isomorphisms \(\rho:G\to G'\) and \(\eta:I\to I'\) that
preserve the attributes selected by \(q\), preserve declared interpretation
relations, and make the mapping commute:
\begin{equation}
  \rho\bigl(\mu(u)\bigr)=\mu'\bigl(\eta(u)\bigr)
  \qquad\text{for every }u\in V_I.
\end{equation}
\end{definition}

\subsection{Provenance}

\begin{definition}[Extensional provenance]
For a derived object \(y\), base-graph provenance is
\begin{equation}
  \pi_G(y)=(V_y,E_y),
  \qquad V_y\subseteq V,\quad E_y\subseteq E.
\end{equation}
For an interpretation node \(u\), provenance agrees with the mapping:
\begin{equation}
  \pi_G(u)=\bigl(V_{\mu(u)},E_{\mu(u)}\bigr).
\end{equation}
\end{definition}

\begin{definition}[Derivation provenance]
Derivation provenance is
\begin{equation}
  \delta(y)=(o,\theta,Y_{\mathrm{parent}}),
\end{equation}
where \(o\) is the producing operator, \(\theta\) is its serialized parameter
state, and \(Y_{\mathrm{parent}}\) is the parent collection declared by the
operator contract.
\end{definition}

Extensional provenance records where an object originates; derivation
provenance records how it was produced. Neither is model attribution, which
assigns relevance rather than structural origin.

\subsection{Operators and programs}

\begin{definition}[Typed operator]
An operator of arity $k$ is a typed function
\begin{equation}
  o:\prod_{j=1}^{k}\mathcal{A}_{C_{\mathrm{in},j}}\to
    \mathcal{A}_{C_{\mathrm{out}}},
\end{equation}
where \(\mathcal{A}_C\) is the set of valid Abstract Graphs satisfying
contract \(C\). A contract declares accepted graph kinds and attributes,
read/write effects on \(G\), \(I\), and \(\mu\), materialisation and coverage,
determinism, relation and provenance semantics, failure conditions, and
complexity.
\end{definition}

The operator family is closed over Abstract Graphs by construction. For
type-compatible unary operators, a sequential program is
\begin{equation}
  P=o_k\circ\cdots\circ o_1,
  \qquad P(A)=o_k\bigl(\cdots o_2(o_1(A))\cdots\bigr).
\end{equation}
Compatibility requires the output contract of \(o_i\) to imply the input
contract of \(o_{i+1}\). Programs may additionally use declared additive,
product, conditional, or bounded iterative composition. A serialized operator
program is denoted by \(\sigma_P(P)\). Closure provides composability and
well-formed outputs; it does not imply refinement of the input representation
or monotone growth in discrimination.

\subsection{Representation invariants}
\label{app:invariants}

The following invariants define valid evidence in this paper.

\begin{enumerate}
  \item \textbf{Mapping validity.} Every node and edge in every \(\mu(u)\)
    belongs to \(G\), with matching semantic attributes.
  \item \textbf{Interpretation totality.} \(\mu(u)\) is defined for every
    \(u\in V_I\).
  \item \textbf{Provenance completeness.} \(\pi_G(y)\) contains all and only
    the base nodes and edges required by the semantics of \(y\), and
    \(\delta(y)\) identifies its derivation when required.
  \item \textbf{Permutation invariance.} For permitted \(\rho\) and
    deterministic \(P\),
    \begin{equation}
      P\bigl(\rho(A)\bigr)\equiv_{q,R_I}\rho\bigl(P(A)\bigr).
    \end{equation}
  \item \textbf{Determinism.} Fixed input, configuration, program, software
    version, and random state produce equivalent outputs and identities.
  \item \textbf{Composition closure.} Adjacent operators have compatible
    contracts, and successful outputs satisfy their output contracts.
  \item \textbf{Identity consistency.} If \(s_1\equiv_q s_2\), their
    canonical certificates and configured feature identities are equal.
  \item \textbf{Serialization round trip.} For every supported serializer
    \(\sigma\) and deserializer \(\sigma^{-1}\),
    \begin{equation}
      \sigma^{-1}\bigl(\sigma(z)\bigr)\equiv z
    \end{equation}
    under the declared equivalence for programs or representation states.
\end{enumerate}
